\documentclass[11pt]{amsart}
\usepackage[T1]{fontenc}
\usepackage{lmodern}
\usepackage{amsmath,amssymb,amsthm,mathtools}
\usepackage{xcolor}
\usepackage[colorlinks=true,linkcolor=blue!55!black,citecolor=blue!55!black,urlcolor=blue!55!black]{hyperref}

\newtheorem{theorem}{Theorem}[section]
\newtheorem{lemma}[theorem]{Lemma}
\newtheorem{proposition}[theorem]{Proposition}
\newtheorem{corollary}[theorem]{Corollary}
\theoremstyle{remark}
\newtheorem{remark}[theorem]{Remark}

\title[Distinct cycle lengths]{The distinct-cycle-length probability is strictly decreasing after 30}
\author{Zijian Zeng}
\address{Institute of Computer Science and Digital Innovation, UCSI University, Kuala Lumpur 56000, Malaysia}
\email{zijianzeng@foxmail.com}
\email{1002266693@ucsiuniversity.edu.my}
\date{August 30, 2026}
\subjclass[2020]{Primary 05A16; Secondary 60C05, 05A17}
\keywords{random permutation, distinct cycle lengths, coefficient monotonicity, infinite product, exact computation}
\hypersetup{
  pdftitle={The distinct-cycle-length probability is strictly decreasing after 30},
  pdfauthor={Zijian Zeng}
}

\begin{document}

\begin{abstract}
Let \(q_n\) be the probability that the cycle lengths of a uniformly random
permutation in \(S_n\) are pairwise distinct. We prove
\(q_{n+1}<q_n\) for every \(n\ge30\). The proof factors the difference
generating function \((1-z)\prod_{k\ge1}(1+z^k/k)\) as \(W(z^2)J(z)\).
An elementary recurrence gives a uniform coefficient estimate for \(W\),
while absolute moment and pointwise bounds for \(J\) preserve a parity
cancellation at \(z=-1\). Explicit constants make every difference negative
from degree 5000 onward. The remaining finite comparisons are certified by
exact integer arithmetic and independently by outward-rounded rational
intervals.
\end{abstract}

\maketitle

\section{Introduction}

For \(n\ge0\), let \(q_n\) denote the probability that no two cycles of a
uniformly random permutation in \(S_n\) have the same length, with \(q_0=1\).
The cycle index gives the ordinary generating function
\begin{equation}\label{eq:Q}
 Q(z):=\sum_{n\ge0}q_nz^n
 =\prod_{k\ge1}\left(1+\frac{z^k}{k}\right).
\end{equation}
Equivalently, \(q_n\) is the sum of \(1/\prod_{a\in A}a\) over sets \(A\)
of distinct positive integers whose sum is \(n\). This is the
distinct-part reciprocal-weight model studied in the classical literature;
see Lehmer~\cite{Lehmer1972}, Greene and Knuth~\cite{GreeneKnuth1981}, and
OEIS A007838~\cite{OEISA007838}.

Greene and Knuth proved
\[
 q_n=e^{-\gamma}\left(1+\frac1n\right)
      +O\!\left(\frac{\log n}{n^2}\right).
\]
Flajolet, Fusy, Gourdon, Panario, and Pouyanne later obtained a full
root-of-unity expansion with logarithmic and periodic terms
\cite[Section~3 and Proposition~1]{FlajoletEtAl2006}. Those qualitative
expansions do not by themselves provide a numerical threshold for the sign
of one-step differences. The problem of proving
\begin{equation}\label{eq:target}
 q_{n+1}<q_n\qquad(n\ge30)
\end{equation}
was recorded as open, together with an exact computation through 5000, in
TheoremDB~\cite{TheoremDBP2540}.

Our proof does not require an effective version of the full asymptotic
expansion. Instead, it isolates the only coefficient tail of order
\(n^{-2}\), keeps the cancellation between the two parity classes, and
bounds everything else in elementary coefficient norms.

\begin{theorem}\label{thm:main}
For every integer \(n\ge30\), the probability that a uniformly random
permutation in \(S_{n+1}\) has pairwise distinct cycle lengths is strictly
smaller than the corresponding probability in \(S_n\); that is,
\[
q_{n+1}<q_n.
\]
\end{theorem}

\section{The difference factorization}

Put
\[
 F(z):=(1-z)Q(z)=\sum_{m\ge0}f_mz^m,
 \qquad f_m=q_m-q_{m-1},
\]
where \(q_{-1}:=0\). Thus Theorem~\ref{thm:main} is equivalent to
\(f_m<0\) for every \(m\ge31\).

Define
\begin{align}
 W(t)&:=\exp\!\left(-\frac12\sum_{k\ge2}\frac{t^k}{k^2}\right),
                                                        \label{eq:W}\\
 R(z)&:=\sum_{r\ge3}\frac{(-1)^{r-1}}r
              \sum_{k\ge2}\frac{z^{rk}}{k^r},          \label{eq:R}\\
 K(z)&:=e^{R(z)},\qquad
 J(z):=(1+z)e^{-z}K(z).                                 \label{eq:J}
\end{align}

\begin{lemma}[Difference factorization]\label{lem:factor}
As an identity of formal power series,
\begin{equation}\label{eq:factor}
 F(z)=W(z^2)J(z).
\end{equation}
Moreover, all series involved are absolutely convergent at the boundary
points used below, and
\begin{equation}\label{eq:boundary}
 W(1)J(1)=e^{-\gamma},\qquad J(-1)=0.
\end{equation}
\end{lemma}

\begin{proof}
Expanding the logarithm in~\eqref{eq:Q}, first separate the factor with
\(k=1\), and then combine the terms linear in \(z^k/k\):
\begin{align*}
 \log F(z)
 &=\log(1-z^2)+\sum_{k\ge2}\log\left(1+\frac{z^k}{k}\right)\\
 &=\log(1+z)-z
   -\frac12\sum_{k\ge2}\frac{z^{2k}}{k^2}+R(z).
\end{align*}
Exponentiation proves~\eqref{eq:factor}. The absolute coefficient bounds in
Sections~\ref{sec:W} and~\ref{sec:J} justify evaluation at \(z=\pm1\).
The factor \(1+z\) gives \(J(-1)=0\). At \(z=1\), the classical Euler
product
\[
 e^{-\gamma}=\prod_{k\ge1}\left(1+\frac1k\right)e^{-1/k}
\]
is exactly the value of the right side of~\eqref{eq:factor}, proving the
first identity in~\eqref{eq:boundary}.
\end{proof}

Write \(J(z)=\sum_{b\ge0}j_bz^b\). Once absolute convergence is known,
Lemma~\ref{lem:factor} gives the crucial parity identity
\begin{equation}\label{eq:parity}
 \sum_{b\equiv m\;({\rm mod}\,2)}j_b
 =\frac{J(1)+(-1)^mJ(-1)}2=\frac{J(1)}2.
\end{equation}

\section{The principal factor}\label{sec:W}

Let
\[
 W(t)=1-\sum_{a\ge2}v_at^a,
 \qquad
 \eta:=W(1)=\exp\!\left(-\frac{\zeta(2)-1}{2}\right).
\]

\begin{lemma}\label{lem:Wsign}
For every \(a\ge2\),
\begin{equation}\label{eq:vbasic}
 0<v_a\le\frac1{2a^2}.
\end{equation}
For every \(a\ge4\),
\begin{equation}\label{eq:Wapprox}
 \left|v_a-\frac{\eta}{2a^2}\right|
 \le\frac{\log a+1}{2a^3}.
\end{equation}
Finally,
\begin{equation}\label{eq:Wnorm}
 \sum_{a\ge0}|[t^a]W(t)|<\frac32.
\end{equation}
\end{lemma}

\begin{proof}
Logarithmic differentiation of~\eqref{eq:W} gives
\begin{equation}\label{eq:vrec}
 a v_a=\frac1{2a}-\frac12\sum_{b=2}^{a-2}\frac{v_b}{a-b}.
\end{equation}
The claim~\eqref{eq:vbasic} is immediate for \(a=2,3\). Suppose it holds
below \(a\). Partial fractions give
\[
 \sum_{b=2}^{a-2}\frac1{b^2(a-b)}
 =\frac1a\sum_{b=2}^{a-2}\frac1{b^2}
  +\frac2{a^2}\sum_{b=2}^{a-2}\frac1b<\frac2a.
\]
Here we used
\(\sum_{b\ge2}b^{-2}<\sum_{b\ge2}[b(b-1)]^{-1}=1\) and
\(\sum_{b=2}^{a-2}b^{-1}<(a-3)/2\). Therefore the right side of
\eqref{eq:vrec} is positive. Dropping its subtracted sum then gives the
upper bound in~\eqref{eq:vbasic}, completing the induction.

Since \(\eta=1-\sum_{b\ge2}v_b\), equation~\eqref{eq:vrec} can be rewritten
as
\[
 v_a=\frac{\eta}{2a^2}
 +\frac1{2a^2}\sum_{b\ge a-1}v_b
 -\frac1{2a^2}\sum_{b=2}^{a-2}\frac{bv_b}{a-b}.
\]
The integral estimate for the tail of \(\sum b^{-2}\), together with
\[
 \sum_{b=2}^{a-2}\frac1{b(a-b)}
 =\frac2a\sum_{b=2}^{a-2}\frac1b\le\frac{2\log a}{a},
\]
now gives
\[
 \left|v_a-\frac{\eta}{2a^2}\right|
 \le\frac{\log a}{2a^3}+\frac1{4a^2(a-2)}.
\]
For \(a\ge4\), the last summand is at most \(1/(2a^3)\), proving
\eqref{eq:Wapprox}. Finally, coefficientwise majorization of the exponential
in~\eqref{eq:W} yields
\[
 \sum|[t^a]W(t)|
 \le\exp\!\left(\frac{\zeta(2)-1}{2}\right)<e^{1/3}<\frac32.
\]
For instance, \(\zeta(2)<5/3\) follows by summing through 5 and bounding the
remaining decreasing tail by \(\int_5^\infty x^{-2}\,dx\).
\end{proof}

\section{Norms and a pointwise bound for the remainder factor}\label{sec:J}

Write \(R(z)=\sum_{m\ge1}\rho_mz^m\) and
\(K(z)=\sum_{m\ge0}\kappa_mz^m\).

\begin{lemma}\label{lem:Rbounds}
The coefficients of \(R\) satisfy
\begin{equation}\label{eq:Rpoint}
 |\rho_m|\le\frac{19}{m^3}\quad(m\ge1),
\end{equation}
and
\begin{equation}\label{eq:Rmoments}
 A_0:=\sum_m|\rho_m|<\frac18,
 \qquad
 A_1:=\sum_m m|\rho_m|\le1.
\end{equation}
\end{lemma}

\begin{proof}
For a term in~\eqref{eq:R} with \(m=rk\), multiplication by \(m^3\)
turns its absolute value into \(r^2/k^{r-3}\). For \(6\le m<100\), exact
rational enumeration of the divisors gives maximum \(113/6\), attained at
\(m=12\). For \(m\ge100\), the term with \(r=3\) is at most 9, the term
with \(r=4\) is at most \(16/25\), and the terms with \(r\ge5\) and
\(k\ge3\) total at most
\[
 \sum_{r\ge5}\frac{r^2}{3^{r-3}}=\frac{31}{6}.
\]
There is at most one term with \(k=2\); it is bounded by
\(50^2/2^{47}<1/100\). The resulting sum is below 15. This proves
\eqref{eq:Rpoint}; the short finite enumeration is independently executable
from the accompanying source.

For~\eqref{eq:Rmoments}, reverse the two nonnegative majorant sums. First,
\[
 A_0\le\sum_{r\ge3}\frac{\zeta(r)-1}{r}
 \le\frac13\sum_{k\ge2}\frac1{k^2(k-1)}
 =\frac{2-\zeta(2)}3<\frac18.
\]
The last inequality follows from
\[
 \zeta(2)>\sum_{k=1}^5\frac1{k^2}
 +\int_6^\infty\frac{dx}{x^2}
 =\frac{5869}{3600}>\frac{13}{8}.
\]
Second,
\[
 A_1\le\sum_{k\ge2}\sum_{r\ge3}\frac1{k^{r-1}}
 =\sum_{k\ge2}\frac1{k(k-1)}=1.
\]
\end{proof}

\begin{lemma}\label{lem:Jbounds}
The coefficients of \(J\) satisfy
\begin{equation}\label{eq:Jmoments}
 \sum_m|j_m|<\frac{16}{7},
 \qquad
 \sum_m m|j_m|<\frac{40}{7},
\end{equation}
and the pointwise estimate
\begin{equation}\label{eq:Jpoint}
 |j_m|<\frac{88}{m^3}\qquad(m\ge2000).
\end{equation}
\end{lemma}

\begin{proof}
Equation~\eqref{eq:Rmoments} and exponential majorization give
\begin{equation}\label{eq:Kmoments}
 \sum_m|\kappa_m|<\frac87,
 \qquad
 \sum_m m|\kappa_m|<\frac87.
\end{equation}
For the pointwise estimate, consider the coefficient of \(R(z)^s\). At
least one member of every ordered \(s\)-tuple of positive indices summing to
\(m\) is at least \(m/s\). Thus Lemma~\ref{lem:Rbounds} gives
\[
 |[z^m]R(z)^s|\le\frac{19s^4A_0^{s-1}}{m^3}.
\]
The exponential generating identity
\[
 \sum_{s\ge1}\frac{s^4a^{s-1}}{s!}
 =e^a(a^3+6a^2+7a+1)
\]
and \(A_0<1/8\) show that this sum is below \(9/4\). For completeness,
\[
 e^{1/8}<1+\frac18+\frac1{128}+\frac1{2976}
 =\frac{13489}{11904}<\frac{17}{15},
\]
where \(k!\ge6\cdot4^{k-3}\) bounds the exponential series from the cubic
term onward. Therefore
\begin{equation}\label{eq:Kpoint}
 |\kappa_m|<\frac{43}{m^3}.
\end{equation}

Let \(B(z)=(1+z)e^{-z}=\sum b_mz^m\). Then \(b_0=1,b_1=0\) and
\[
 b_m=\frac{(-1)^m(1-m)}{m!}\quad(m\ge2),
\]
so \(\sum|b_m|=2\) and \(\sum m|b_m|=e<3\). Convolution with
\eqref{eq:Kmoments} proves~\eqref{eq:Jmoments}.

For \(m\ge2000\), split \(j_m=\sum_i b_i\kappa_{m-i}\) at \(i=m/2\).
Convexity gives \((1-i/m)^{-3}\le1+14i/m\) on \([0,1/2]\), so
\eqref{eq:Kpoint} bounds the first part by
\[
 \frac{43}{m^3}\left(2+\frac{42}{m}\right)<\frac{87}{m^3}.
\]
The remaining part is less than \(m^{-3}\): if
\(s=\lfloor m/2\rfloor\), its absolute value is at most
\((8/7)s/(s+1)!<(8/7)/s!\), and
\(s!\ge2^{s-1}>(216/7)s^3> (8/7)m^3\) for \(s\ge1000\).
This proves~\eqref{eq:Jpoint}.
\end{proof}

\section{The uniform tail}

\begin{proposition}\label{prop:tail}
For every integer \(m\ge5000\), \(f_m<0\).
\end{proposition}

\begin{proof}
Write \(w_a=[t^a]W(t)\), so that \(w_0=1,w_1=0\), and
\(w_a=-v_a\) for \(a\ge2\). From~\eqref{eq:factor},
\begin{equation}\label{eq:conv}
 f_m=\sum_{\substack{0\le b\le m\\b\equiv m\;({\rm mod}\,2)}}
 j_bw_{(m-b)/2}.
\end{equation}
Split the sum at \(b=m/2\). In the first half put \(x=m-b\); then
\(x\ge m/2\), and Lemma~\ref{lem:Wsign} replaces \(w_{x/2}\) by
\(-2\eta/x^2\). Extend the resulting constant main term across the full
parity sum. Equations~\eqref{eq:boundary} and~\eqref{eq:parity} make it
exactly
\begin{equation}\label{eq:mainterm}
 -\frac{2\eta}{m^2}\frac{J(1)}2=-\frac{e^{-\gamma}}{m^2}.
\end{equation}

We bound the four errors explicitly. First, the derivative of
\(-2\eta/x^2\) has absolute value at most \(32/m^3\) on
\([m/2,m]\). The first moment in~\eqref{eq:Jmoments} therefore bounds its
variation by \(1280/(7m^3)\). Second,~\eqref{eq:Wapprox}, \(x\ge m/2\),
and the zeroth moment in~\eqref{eq:Jmoments} bound the
coefficient-approximation error by
\(512(\log m+1)/(7m^3)\). Third, in the half-tail \(b>m/2\),
equation~\eqref{eq:Jpoint} gives \(|j_b|<704/m^3\); summing the matching
\(W\)-coefficients with~\eqref{eq:Wnorm} costs less than \(1056/m^3\).
Finally, \(\sum_{b>m/2}|j_b|<176/m^2\), so extending the main parity sum
across that tail costs less than \(352/m^4<m^{-3}\). Combining these
estimates with~\eqref{eq:mainterm},
\begin{equation}\label{eq:tailbound}
 f_m\le-\frac{e^{-\gamma}}{m^2}
 +\frac1{m^3}\left[
 \frac{1280}{7}+1057+\frac{512}{7}(\log m+1)\right].
\end{equation}

It remains to check the explicit endpoint. Since \(H_n-\log n\) decreases
to \(\gamma\), the positive series for the inverse hyperbolic tangent give
\[
 \log2>\frac{56}{81}>\frac23,
 \qquad
 \log3>\frac{263}{240},
\]
and hence
\[
 \gamma<\frac{49}{20}-\log6<\frac23.
\]
Thus \(e^{-\gamma}>1/2\). Also \(\log5000<9\), because
\(e>8/3\) and \((8/3)^9>5000\). At \(m=5000\), the bracket in
\eqref{eq:tailbound} is therefore below
\[
 \frac{13799}{7}<2500=\frac m2.
\]
The derivative of the bracket is \(512/(7m)<1/2\) throughout
\(m\ge5000\), so the strict comparison persists for all larger \(m\).
Equation~\eqref{eq:tailbound} is consequently negative.
\end{proof}

\section{The finite certificate and completion of the proof}

For the finite range we used two independent rigorous representations.
First set \(D=5000!\) and expand~\eqref{eq:Q} with descending updates. Every
scaled coefficient \(c_m=Dq_m\) is an integer: a summand indexed by a set
\(A\) of distinct parts has denominator \(\prod_{a\in A}a\), and
if \(\sum_{a\in A}a=m\le5000\), then
\[
 \frac{m!}{\prod_{a\in A}a}
 =\frac{m!}{\prod_{a\in A}a!}\prod_{a\in A}(a-1)!
\]
is an integer, since the first factor on the right is a multinomial
coefficient.  As \(m!\) divides \(D\), it follows that
\(D/\prod_{a\in A}a\) is an integer as well.  Equivalently, for each part
\(k\), the program applies the descending exact update
\(c_d\leftarrow c_d+c_{d-k}/k\); it rejects the run if any displayed division
has a nonzero remainder. Exact integer comparison gives
\begin{equation}\label{eq:finite}
 c_m<c_{m-1}\qquad(31\le m\le5000).
\end{equation}

Independently, a fixed-point checker encloses every coefficient of each
partial product between integer multiples of \(10^{-30}\), rounding every
division outward. It accepts~\eqref{eq:finite} only if the upper endpoint for
\(q_m\) lies strictly below the lower endpoint for \(q_{m-1}\). A second run
at scale \(10^{-40}\) gives the same comparisons. At scale \(10^{-30}\),
the smallest separating gap is
\[
 22282582258338445775491\times10^{-30},
\]
at \(m=4998\), while the largest interval width is only
\(12408\times10^{-30}\). These integers are diagnostic summaries; the
verifier recomputes every enclosure from the product.

\begin{proof}[Proof of Theorem~\ref{thm:main}]
The exact certificate proves \(f_m<0\) for \(31\le m\le5000\), and
Proposition~\ref{prop:tail} proves it for every \(m\ge5000\). Hence
\(f_m=q_m-q_{m-1}<0\) for all \(m\ge31\). Taking \(m=n+1\) proves
\(q_{n+1}<q_n\) for every \(n\ge30\).
\end{proof}

\section*{Data, code, and computational disclosure}

The accompanying source contains the exact common-denominator program, the
independent outward-rounded interval checker, an exact rational checker for
the constants in Lemmas~\ref{lem:Rbounds}--\ref{lem:Jbounds}, and an exact
low-degree reconstruction of~\eqref{eq:factor}. All use only the Python 3
standard library. A fast command runs the interval, constant, and structural
checks; a full command additionally repeats the \(5000!\)-denominator
expansion.

Computer algebra and language-model assistance were used during proof
discovery, adversarial checking, code review, and manuscript preparation.
All theorem statements, derivations, citations, and certificate logic were
subsequently checked by the author. No proof assistant was used.

\bibliographystyle{amsplain}
\bibliography{references}

\end{document}
